%======================================================================
\section{Discussion}\label{sec:discussion}
%======================================================================

\paragraph{What has been settled.} For uniform internally heated convection at infinite Prandtl number, the
fraction of heat leaving through the base obeys $\overline{F_B}\ge c(R\mathcal T\ln R\mathcal T)^{-1/3}$
between no-slip plates and $\overline{F_B}\ge c(R\mathcal T)^{-1/2}$ between stress-free plates, for every
classical solution. The maximum temperature $\mathcal T$ is bounded by $C(1+\ln|\Omega|)$ for every incompressible flow. At
fixed aspect ratio this improves the best unconditional exponents from $2/3$ to $1/3$ (up to a logarithm) and from
$40/29$ to $1/2$. The bound is explicit: $\mathfrak B[\frac12R\mathcal T\Phi]$ with a computable wall function. The
rate $R^{-1/3}$ had appeared before only under the unproved profile hypothesis of \citet{ChenEtAl2024}; the finite-window grid tests do not establish that hypothesis. The numerical values of the constants rest on two
computed norms, which are not interval-certified (Table~\ref{tab:status}).

\paragraph{Why it works.} Bounds from the background method control $\langle wT\rangle$ through quadratic forms
and have to pay for the stable bottom layer through the minimum principle, imposed as a constraint. Here the
minimum principle does the main work directly: since $T\ge0$, the averaged flux is at most the maximal speed times
the mean temperature, which is linear in the temperature. The no-slip condition then makes the bound on the
near-wall speed quadratically small in the height, up to a logarithm, so that its integral, which sets the
integrating factor, grows only like $z^3\ln(1/z)$.

\paragraph{What remains open.}
\begin{enumerate}
\item \emph{Attainment.} The bound is not approached by any flow we computed. Time-dependent simulations and
$F_B$-minimising steady rolls both give $F_B\approx1.3R^{-0.09}$ over the computed ranges; the proof's extremal configuration requires the convective flux near the
base to be close to $z-F_B$, whereas real flows leave the lower third of the layer to conduction. Whether
\emph{some} solution has $F_B\to0$ as fast as a power $R^{-\alpha}$ with $\alpha$ near $1/3$ is unknown. A
heuristic suggests why. Plumes that reach the base still carry roughly the interior temperature, so conduction into
the base stays comparable to the heat generated in the stable layer. On this picture $F_B$ decays only as fast as
that layer thins, much more slowly than our bound.
\item \emph{The logarithm.} It enters only through the $L^\infty$ kernel estimate. Removing it would need
information about the temperature at every height between the wall and the interior, in the spirit of the bilinear
estimate of \citet{ChanilloMalchiodi2025}.
\item \emph{The aspect ratio.} The factor $\ln|\Omega|$ comes from the temperature bound. Whether $\mathcal T$ is
bounded uniformly in the aspect ratio is a question about exit times of diffusions with incompressible drift in a
slab \citep{IyerEtAl2010}. Any uniform bound would make Corollary~\ref{cor:R} uniform as well.
\item \emph{Finite Prandtl number.} The argument requires a pointwise, instantaneous near-wall velocity bound, which
the inertial equations do not supply. Proving the \citet{ChenEtAl2024} hypothesis, at any Prandtl number, would give
$F_B\gtrsim R^{-1/3}$ without the logarithm or the temperature factor.
\item \emph{Rayleigh--B\'enard construction.} At infinite Prandtl number the largest Nusselt maxima found among
two-dimensional mirror-symmetric, layer-spanning steady rolls follow $\mathrm{Nu}\approx0.15\,\mathrm{Ra}^{0.298}$
for $10^6\le\mathrm{Ra}\le10^8$, consistent with the $\mathrm{Ra}^{3/10}$ high-wavenumber asymptotics, and the
wide-roll maximum that gives $\mathrm{Ra}^{1/3}$ at $\mathrm{Pr}=1$ is absent from the scanned period range above
$\mathrm{Ra}\approx1.8\times10^5$. These results motivate the investigation of three-dimensional,
time-dependent or multi-scale constructions of solutions attaining the upper bound
$\mathrm{Nu}\lesssim\mathrm{Ra}^{1/3}$. Three-dimensional multi-scale steady solutions are known at $\mathrm{Pr}=1$, but up to
$\mathrm{Ra}\sim10^7$ they give $\mathrm{Nu}\sim\mathrm{Ra}^{0.31}$ \citep{MotokiEtAl2021}. A rigorous version
of the $\mathrm{Ra}^{3/10}$ family, via matched asymptotics turned into an existence proof, would be a natural
first target. It would exhibit rigorously a family of solutions whose Nusselt number grows as a power of Ra,
though with exponent $3/10$ rather than $1/3$.
\end{enumerate}

\paragraph{Limitations of this study.} The bounds are proved for classical solutions in two and three dimensions;
the computations are two-dimensional. The numerical constants $C_{\rm ns},C_{\rm sf}$ come from grid maximisation
and quadrature, with numerical sensitivity checks, not uniform error enclosures and a second, more conservative route through Carlson majorants, but not from
interval arithmetic. Their finiteness is proved, and the constants of Theorem~\ref{thm:main} depend on them only
weakly. The simulations are four runs at two Rayleigh numbers; three of them are intermittent, and their window
averages have the descriptive sub-block spread shown in Table~\ref{tab:dns}, not a justified uncertainty for long-time statistics. The maximum temperatures observed over finite
windows are not bounds on the long-time maximum. The
steady-roll results rest on finitely many Rayleigh numbers, the largest $10^8$ (Rayleigh--B\'enard) and
$2.8\times10^7$ (internally heated), and on mirror-symmetric, layer-spanning rolls. The continued roll branch was
sampled only at six Rayleigh numbers between $10^5$ and $10^7$ (Rayleigh--B\'enard) and at $R=1.6\times10^6$
(internally heated), and only over the period ranges stated; other solution families, such as rolls stacked in
the vertical, were not explored.

\paragraph{Authorship and data.} This work was produced by directed AI agents for Evidence Press. Code, data and
logs sufficient to regenerate every figure, table and number are in the accompanying bundle.
